\SourceLeaf{095}
III, IV, V, VII ad hunc finem combinare liceat; attamen levis attentio docet, \textit{duas} sufficere, alteram ex illis, alteram ex his, quum reliquae in his atque prioribus aequationibus conditionalibus iam contentae esse debeant. Aequatio itaque conditionalis sexta nobis erit
\[
\begin{aligned}
&\log\sin(v^{(3)}-v^{(2)}-0''067)-\log\sin(v^{(5)}-v^{(4)}-0''067)\\
&+\log\sin(v^{(14)}-v^{(17)}-0''640)-\log\sin(v^{(2)}-v^{(0)}-0''640)\\
&+\log\sin(v^{(6)}-v^{(5)}-0''107)-\log\sin(v^{(17)}-v^{(16)}-0''107)=0
\end{aligned}
\]
atque septima
\[
\begin{aligned}
&\log\sin(v^{(2)}-v^{(1)}-0''419)-\log\sin(v^{(12)}-v^{(11)}-0''419)\\
&+\log\sin(v^{(14)}-v^{(17)}-0''640)-\log\sin(v^{(2)}-v^{(0)}-0''640)\\
&+\log\sin(v^{(13)}-v^{(11)}-0''432)-\log\sin(v^{(17)}-v^{(15)}-0''432)=0
\end{aligned}
\]
quibus respondent aequationes complexus (13)
\[
\begin{aligned}
+25={}&+4{,}31(0)-153{,}88(2)+149{,}57(3)+39{,}11(4)-79{,}64(5)\\
&+40{,}53(6)+31{,}90(14)+275{,}39(16)-307{,}29(17)\\
-3={}&+4{,}31(0)-24{,}16(1)+19{,}85(2)+36{,}11(11)-28{,}59(12)\\
&-7{,}52(13)+31{,}90(14)+29{,}06(15)-60{,}96(17)
\end{aligned}
\]

Quodsi iam singulis directionibus eandem certitudinem tribuimus, statuendo \(p^{(0)}=p^{(1)}=p^{(2)}\ \text{etc.}=1\), correlataque septem aequationum conditionalium, eo ordine, quem hic sequuti sumus, per \(A,B,C,D,E,F,G\) denotamus, horum determinatio petenda erit ex aequationibus sequentibus:
\[
\begin{aligned}
-1{,}368={}&+6A-2B-2C-2D+184{,}72F-19{,}85G\\
+1{,}773={}&-2A+6B+2C+2E-153{,}88F-20{,}69G\\
+1{,}042={}&-2A+2B+6C-2D-2E+181{,}00F+108{,}40G\\
-0{,}813={}&-2A-2C+6D+2E-462{,}51F-60{,}96G\\
-0{,}750={}&+2B-2C+2D+6E-307{,}29F-133{,}65G\\
+25={}&+184{,}72A-153{,}88B+181{,}00C-462{,}51D-307{,}29E\\
&+224868F+16694{,}1G\\
-3={}&-19{,}85A-20{,}69B+108{,}40C-60{,}96D-133{,}65E\\
&+16694{,}1F+8752{,}39G
\end{aligned}
\]

Hinc deducimus per eliminationem
\SourceLeaf{096}
\[
\begin{aligned}
A&=-0{,}225\\
B&=+0{,}344\\
C&=-0{,}088\\
D&=-0{,}171\\
E&=-0{,}323\\
F&=+0{,}000215915\\
G&=-0{,}00547462
\end{aligned}
\]

Iam errores maxime plausibiles habentur per formulas:
\[
\begin{aligned}
(0)&=-C+4{,}31F+4{,}31G\\
(1)&=-B-24{,}16G\\
(2)&=-A+B+C-153{,}88F+19{,}85G
\end{aligned}
\]
etc., unde prodeunt valores numerici
\[
\begin{array}{r@{\;}l@{\qquad}|@{\qquad}r@{\;}l}
(0)=&+0''065&(9)=&+0''021\\
(1)=&-0{,}212&(10)=&+0{,}054\\
(2)=&+0{,}339&(11)=&-0{,}219\\
(3)=&-0{,}193&(12)=&+0{,}501\\
(4)=&+0{,}233&(13)=&-0{,}282\\
(5)=&-0{,}071&(14)=&-0{,}256\\
(6)=&-0{,}162&(15)=&+0{,}164\\
(7)=&-0{,}481&(16)=&+0{,}230\\
(8)=&+0{,}406&(17)=&-0{,}139
\end{array}
\]

Summa quadratorum horum errorum invenitur \(=1{,}2288\); hinc error medius unius directionis, quatenus e 18 directionibus observatis erui potest,
\[
=\BareRoot{\frac{1{,}2288}{7}}=0''4190
\]

\Article{25}
Ut etiam pars altera theoriae nostrae exemplo illustretur, indagamus praecisionem, qua latus Falkenberg-Breithorn e latere Wilsede-Wulfsode adiumento observationum compensatarum determinatur. Functio \(u\), per quam illud in hoc casu exprimitur, est
